\section{Sensor Characterization}
\label{appendix:sensor_characterization}


\subsection{VL53L0X Time-of-Flight Sensor}

The position sensor is one of the defining sources of difficulty on this
platform, so we characterize its behavior carefully before turning to the
controllers. Ball position is read by two VL53L0X (STMicroelectronics)
Time-of-Flight sensors, each with a programmable timing budget that sets the
integration time per measurement. We mount one at each end of the nearly
30\,cm arc (radius 2.101\,m, approximately linear) so that whichever half of
the arc the ball occupies is always read by the nearer sensor: the left sensor
covers the left half, the right sensor the right half. This nominal span
corresponds to $\pm 0.071$\,rad at $R = 2.101$\,m; logged ball angles
occasionally reach the $\pm 0.081$\,rad empirical limit used elsewhere,
because each sensor's 0--170\,mm range can report the ball marginally past the
nominal arc ends. Both share a single I$^2$C bus,
which precludes reading them simultaneously without interference.

\subsection{Sequential Mode}

In sequential mode we trigger and read the two sensors alternately. The
total cycle time we measure follows
$T_\text{cycle} \approx 2\times\text{budget} + 12\,\text{ms}$,
where the $\sim$12\,ms overhead comes from the I$^2$C transactions and the
inter-trigger delay. To map the tradeoff between noise and update rate, we
held the ball static at three positions, the center of the arc and each
extreme, and averaged 50 samples per position at four timing budgets
(Table~\ref{tab:seq_char}). The minimum programmable budget is 20\,ms, which
caps the sequential update rate at ${\approx}18.6$\,Hz.

\begin{table*}[!t]
\centering
\caption{VL53L0X sequential mode characterization. All values are
         mean\,$\pm\,\sigma$ (mm), averaged over 50 static measurements.
         $^\dagger$Interpolated by applying the $\sigma$ scaling factor
         observed across measured entries; means are physically stable.}
\label{tab:seq_char}
\resizebox{\textwidth}{!}{%
\begin{tabular}{ccc cc cc cc}
\toprule
\multirow{2}{*}{\shortstack{\textbf{Budget}\\(ms)}}
  & \multirow{2}{*}{\shortstack{\textbf{Cycle}\\(ms)}}
  & \multirow{2}{*}{\shortstack{\textbf{Rate}\\(Hz)}}
  & \multicolumn{2}{c}{\textbf{Center}}
  & \multicolumn{2}{c}{\textbf{Left Extreme}}
  & \multicolumn{2}{c}{\textbf{Right Extreme}} \\
\cmidrule(lr){4-5}\cmidrule(lr){6-7}\cmidrule(lr){8-9}
 & & & $L$ & $R$ & $L_\text{min}$ & $R_\text{max}$
       & $L_\text{max}$ & $R_\text{min}$ \\
\midrule
100  & 211.7 & 4.72  
  & $161.80\pm0.85$ & $175.76\pm1.06$
  & $48.32\pm0.99$         & $210.46\pm1.01$
  & $181.22\pm0.97$        & $40.32\pm1.00$ \\
50   & 113.2 & 8.84  
  & $161.34\pm1.24$ & $175.46\pm1.51$
  & $48.20\pm1.44^\dagger$ & $210.30\pm1.47^\dagger$
  & $180.68\pm1.54$        & $40.58\pm1.14$ \\
34.5 & 82.1  & 12.17 
  & $160.54\pm1.62$ & $175.64\pm1.68$
  & $48.08\pm1.90^\dagger$ & $210.12\pm1.93^\dagger$
  & $180.44\pm2.11$        & $41.62\pm1.78$ \\
20   & 53.8  & 18.59 
  & $160.22\pm2.83$ & $176.42\pm2.81$
  & $47.78\pm3.29^\dagger$ & $209.74\pm3.36^\dagger$
  & $180.12\pm3.22$        & $42.34\pm3.01$ \\
\bottomrule
\end{tabular}%
}
\end{table*}

Two systematic effects appear consistently across all configurations:

\begin{enumerate}

\item \textbf{Near-field blind zone:} The VL53L0X datasheet~\cite{vl53l0x_datasheet}
(\S6) specifies a minimum measurable distance dependent on target
reflectance and ambient IR levels.
Below this threshold ($\approx$30--50\,mm in our configuration), the
return photon pulse arrives before the SPAD timing gate opens, causing
one of two failure modes:
\begin{itemize}
  \item The sensor reports \texttt{RangeStatus}~$= 4$ (signal fail),
        which the firmware outputs as $-1$ and the host-side parser
        rejects, the previous valid reading is held.
  \item The sensor returns an erroneously large distance (typically
        $\approx$40--48\,mm regardless of true position), which passes
        the host-side validity check ($<300$\,mm) and is accepted as
        a legitimate measurement.
\end{itemize}

\textbf{Effect on ball-angle estimation:}
As the ball approaches the arc end nearest a sensor, the measured
ball-to-sensor distance shrinks toward zero; near the edge it falls
below the near-field threshold and the reading is no longer valid.
The blind-zone onset occurs at approximately
$\theta_\text{blind} \approx 0.071$\,rad.
Beyond this angle, the reported position \emph{saturates}: the
controller observes the ball at approximately 0.071\,rad, and when
the ball is stationary or moving very slowly, the velocity estimate
appears near zero, even as the true ball angle increases to 0.08\,rad.
If the ball is moving through the blind zone at speed, the saturation
duration is brief and typically inconsequential because the next
reading captures the ball at a different position. When the ball
reverses direction and exits the blind zone, normal tracking resumes.

\item \textbf{Inter-sensor bias:} On this characterization unit the right
sensor reads ${\approx}14$--$16$\,mm higher than the left at center despite
symmetric physical placement. This is an intrinsic per-unit offset rather
than a physical asymmetry: after a later sensor-board swap the deployed board
carries the opposite bias, with the left sensor reading ${\approx}30$\,mm
higher (Sec.~\ref{appendix:sensor_offset}). Either way the offset must be
calibrated out so the two halves of the arc report a consistent zero.

\end{enumerate}

Measurement noise scales inversely with the timing budget, roughly tripling
as the budget drops from 100\,ms to 20\,ms, so faster sensing comes directly
at the cost of noisier readings.

\subsection{Continuous Mode}

Continuous mode ranges both sensors in parallel rather than alternately,
so each unit runs on its own ${\sim}33$\,ms budget (nominally ${\sim}30$\,Hz)
instead of the two sharing the ${\sim}82$\,ms sequential cycle (12\,Hz).
We rejected it for two reasons. The first is optical crosstalk between the co-located
sensors: the VL53L0X uses a 940\,nm VCSEL infrared
emitter~\cite{vl53l0x_datasheet}, and when both emit at once, the signal from
one is picked up by the other's receiver and corrupts the reading. The second
is noise: even setting crosstalk aside, levels run roughly 5 to 30 times
larger than in sequential mode at the same budget (Table~\ref{tab:cont_char}),
enough to render position estimates unusable for control. We therefore use sequential
mode throughout.

\begin{table}[!t]
\centering
\caption{VL53L0X continuous mode at center position
         (mean\,$\pm\,\sigma$, mm), 50 samples. Update rates were
         not reliably characterized in this mode due to crosstalk
         effects. $^\ddagger$Partial outlier
         exclusion applied; fault rate 3--5\%.}
\label{tab:cont_char}
\begin{tabular}{ccc}
\toprule
\textbf{Budget (ms)} & $L$ (mm) & $R$ (mm) \\
\midrule
50   & $217.5\pm38.9$          & $117.2\pm10.0$ \\
34.5 & $179.1\pm29.8^\ddagger$ & $144.0\pm8.9^\ddagger$ \\
20   & $147.1\pm41.2$          & $170.9\pm38.1$ \\
\bottomrule
\end{tabular}
\end{table}

\subsection{Simulation Sensor Model}
\label{appendix:sensor_sim_model}

To narrow the sim-to-real gap, our training environment
(\texttt{BalancerSim}) reproduces three of these sensor effects:

\begin{enumerate}
  \item \textbf{Measurement staleness:} Ball position updates at
        $\sim$10\,Hz (geometric inter-arrival, mean $= 2$ control steps).
        Between updates the controller observes the last valid reading
        (this matches the deployed effective ball-update rate quantified
        in the selected-configuration subsection below).

  \item \textbf{Observation noise:} Additive Gaussian noise with
        $\sigma_\theta \approx 0.0005$--$0.002$\,rad
        ($\approx 1.0$--$4.2$\,mm of arc, position) and
        $\sigma_{\dot\theta} \approx 0.02$--$0.04$\,rad/s
        ($\approx 42$--$84$\,mm/s, velocity, amplified by
        finite-difference estimation).
        In domain-randomization mode, $\sigma$ is sampled uniformly
        each episode from the measured hardware range.

  \item \textbf{Near-field blind zone (optional):}
        When $|\theta_\text{true}| > \theta_\text{blind}$ (nominally
        0.064\,rad from the sensor geometry; the calibrated measured
        onset is 0.071\,rad, and both lie within the
        $[0.060, 0.080]$\,rad range randomized under DR), the
        ball-position observation is clamped to
        $\pm\theta_\text{blind}$ and the velocity observation is set
        to $\approx 0$.
        This reproduces the saturation behavior described above:
        the controller sees the ball ``stuck'' at the blind-zone boundary
        while the true state continues evolving.
        The feature is disabled by default
        (\texttt{sensor\_blind\_zone=False}) and must be explicitly
        enabled for training with blind-zone awareness. When combined
        with domain randomization, the threshold is randomized over
        $[0.060, 0.080]$\,rad.
\end{enumerate}

\subsection{Selected Configuration and Control Implications}

We use the default 34.5\,ms timing budget in sequential mode for all
experiments, which yields a 12.17\,Hz sensor update rate (82.1\,ms cycle)
with ${\approx}1.6$--$1.7$\,mm noise at center. This budget is held fixed across
all experiments so that observed performance differences
reflect controller variation alone, not changes in sensor characteristics.
Longer budgets offer marginal noise reduction at a disproportionate rate
penalty; shorter budgets increase noise without sufficient rate benefit for
the tested frequency range. All sensor readings are used raw without
filtering to benchmark controller robustness under realistic sensing
conditions. The 82.1\,ms cycle introduces position-measurement latency that
exceeds the control period (50\,ms at 20\,Hz); all
controllers therefore act on position data that may be up to one full
sensor cycle old. In deployment the effective refresh is slower still:
sampled by the 50\,ms control loop, the fused ball reading updates every
one to two control cycles across the benchmark trials (median 100\,ms,
mean 87\,ms, i.e.\ ${\sim}$10--12\,Hz effective). This is the rate the
simulator's ${\sim}$10\,Hz ball-staleness model reproduces, so the
training environment matches the deployed sensing rate rather than the
faster 12\,Hz rate measured statically on the bench.

\subsection{Calibration Offset}
\label{appendix:sensor_offset}

Post-hoc analysis revealed a systematic offset between the sensor-reported
zero and the true physical center of the arc.  Using the IR beam-break
sensor (which fires when the ball crosses the arc's lowest point) as
ground truth, we measured the ToF sensor readings with the ball
stationary at physical center (5000 samples, dip\,=\,1 for all):

\begin{center}
\begin{tabular}{lcc}
\toprule
 & Left sensor & Right sensor \\
\midrule
Measured at physical center & $193 \pm 4$\,mm & $160 \pm 5$\,mm \\
Calibration used in code    & 175\,mm         & 170\,mm         \\
Error                       & $+18$\,mm       & $-10$\,mm       \\
\bottomrule
\end{tabular}
\end{center}

\noindent
With the earlier code calibration (175/170), the \texttt{process\_dual\_sensors}
routine reports $\theta \approx +4.8$\,mrad when the ball is at physical
center (right sensor path: $-(160-170)/1000/2.101$), and across the
175/170-era trial logs the median ball angle at dip\,=\,1 is $+7.14$\,mrad.
Recalibrating the center constants to the measured equilibrium removes this
residual: in the deployed datasets (\texttt{arcball\_cont\_1M\_calib196\_160},
center constants 189/158) the median ball angle at dip\,=\,1 is
$\approx 0$\,mrad.

\textbf{Noise characterization:}  With the ball stationary at center, the
left sensor spans 187--201\,mm and the right sensor spans
151--168\,mm.  The per-sample standard deviation is
approximately 3\,mm for the left sensor and up to 6\,mm for the right,
the larger figure inflated by intermittent near-field fault readings;
excluding those faults the robust spread is close to 3\,mm for both.
Converted to angle, the typical per-reading noise is
$\sigma_\theta \approx 3\,\text{mm}\, / \,2101\,\text{mm}
\approx 1.4$\,mrad.
The $\pm 4$\,mm and $\pm 5$\,mm figures in the calibration table are the
per-sample spreads over these 5000 stationary, dip-gated samples.

Because \texttt{process\_dual\_sensors} fuses the two units, the ball-angle
estimate that controllers actually receive is tighter than either raw
channel. Over the same stationary samples (from the deployed calibration,
\texttt{arcball\_cont\_1M\_calib196\_160}) its standard deviation is
$\sigma_\theta \approx 1.6$\,mrad ($\approx 3.3$\,mm of arc), with a
peak-to-peak spread near 13\,mm. This is the figure quoted in the main
paper; the 14--17\,mm per-unit spans above are the raw single-sensor
peak-to-peak, so the two are consistent views of the same measurement.
Either way, the deployed-board noise is roughly double the pre-swap
figure and a substantial fraction of the $\pm$21\,mm success tolerance.

\textbf{Note on Table~\ref{tab:seq_char} discrepancy:}
The characterization table reports $L{=}161.8$, $R{=}175.8$ at center,
which differs from the current $L{=}193$, $R{=}160$.  This data was
collected prior to a hardware modification (sensor board swap)
that changed the per-unit offsets.  The post-modification readings
at true physical center are $L\approx 193$, $R\approx 160$.

\textbf{Impact on experiments:}
For the final benchmark (Experiment~1, Table~II in the main paper),
the sensor offset is corrected by recalibrating the center reading:
rather than subtracting a bias from raw readings, the
\texttt{process\_dual\_sensors} routine was updated so that the
center calibration constants produce $\theta = 0$ when the ball is
at physical center.  After recalibration, the left and right baseline
values are no longer 175 and 170; they reflect the new reference point
such that center readings are $\theta \approx 0$ by default.
The offset discussion above applies to earlier ablation runs and
datasets collected before the correction was applied.  For those runs
the bias is constant and smaller than the settling band ($\pm 10$\,mrad);
its main effect is to lengthen settling times, though it also lowers
success rate to a modest degree (for PD+WO, correcting the offset raises
hardware SR from 88\% to 96\%, Appendix~\ref{appendix:implementation}),
which is why it is corrected for the deployed benchmark.
Because these per-unit offsets drift between sessions, recalibrating the
center constants before each experiment campaign is necessary rather than
optional; the specific constants quoted here are representative, and the
exact deployed values vary from run to run. The recalibration procedure
is documented in the reproducibility guide.
